\documentclass[twocolumn]{article}
\usepackage{amsmath,amsfonts,amssymb}
\usepackage{graphicx}
\usepackage{hyperref}
\title{GET-ROPR: Gradient Eligibility Traces for Recurrent Off-Policy RL}
\author{Auto-generated for Assignment CA7}
\begin{document}
\maketitle
\begin{abstract}
We introduce GET-ROPR, an off-policy method that integrates backward-view
TD($\lambda$) eligibility traces into recurrent critics for partially observable
environments. GET-ROPR combines IS-aware trace accumulation, trace clipping,
and hidden-state recomputation to enable long-horizon credit assignment while
maintaining stability with replay buffers.
\end{abstract}

\section{Introduction}
Partial observability necessitates memory; recurrent architectures (RNN/GRU/LSTM)
are commonly used. Off-policy replay introduces stale hidden states and truncation
of credit via limited BPTT. We propose an eligibility-trace based approach that
operates on replayed sequences to recover long-horizon credit while controlling
variance via clipped importance sampling.

In practice, recurrent off-policy agents face two core challenges: (1) hidden-state
staleness — stored hidden states in replay buffers become misaligned as function
approximators change — and (2) limited temporal credit assignment due to
truncated backpropagation through time. GET-ROPR addresses these via three
engineering components: hidden-state recomputation, IS-aware backward-view
eligibility traces (or equivalently forward-view lambda-returns), and conservative
clipping/regularization to stabilize learning.

Contributions: this assignment provides (a) a clear mathematical formulation of
IS-weighted lambda-returns for recurrent critics, (b) a compact PyTorch-style
reference implementation with unit tests and a demo training script, and (c)
practical experimental protocols and ablations to evaluate the effectiveness of
long-horizon traces in partially observable control tasks.

\section{Related work}
GET-ROPR builds on classical temporal-difference methods (TD($\lambda$)), the
literature on off-policy corrections (importance sampling, per-decision IS), and
recent practical recurrent off-policy baselines (e.g., RESeL). We also draw on
ideas from eligibility traces in function approximation literature and modern
actor-critic techniques (SAC/TD3) for practical tuning and stabilization.

Prior works addressing memory and partial observability include recurrent
policy/value networks for MRPs and various strategies for hidden-state handling
during replay (store short prefixes, recompute, or use target refresh). On the
variance reduction side, work on truncated IS products, V-trace, and Retrace
inspire our choice of clipped ratios and controlled trace length.

\section{Background}
We adopt standard notation: observations \(o_t\), actions \(a_t\), rewards \(r_t\),
discounts \(\gamma\), done flags \(d_t\). The recurrent critic \(Q_\phi(o_t,h_t,a_t)\)
produces per-step action-values and evolves hidden states via \(h_{t+1}=f_\phi(o_t,a_t,h_t)\).

\section{Method}
\subsection{Backward-view traces with IS}
Define TD error
\[
\delta_t = r_t + \gamma(1-d_t)Q_{\bar\phi}(o_{t+1},h_{t+1},a_{t+1}) - Q_\phi(o_t,h_t,a_t).
\]
Backward-view clipped IS eligibility trace:
\[
e_t = \bar{\rho}_t\big(\lambda\gamma(1-d_{t-1})e_{t-1} + \nabla_\phi Q_\phi(o_t,h_t,a_t)\big),
\]
where \(\bar{\rho}_t=\min(\rho_t,c_\rho)\) and \(\rho_t=\frac{\pi(a_t|o_t,h_t)}{\mu(a_t|o_t,h_t)}\).
The update direction is \(g=\sum_t \delta_t e_t\).

\subsection{Forward-view surrogate (lambda-returns)}
We use the numerically-stable forward-view surrogate
\[
\mathcal{L}=\tfrac12\sum_t (G_t^\lambda - Q_t)^2,\qquad
G_t^\lambda = Q_t + \sum_{n\ge1} (\gamma\lambda)^{n-1}\Big(\prod_{i=1}^{n-1}\bar{\rho}_{t+i}\Big)\delta_{t+n-1},
\]
which yields gradients equivalent to the backward view under exact arithmetic.

\subsection{Practical considerations and numerical stability}
To ensure stable training we recommend:
\begin{itemize}
  \item Clipping IS ratios \(\bar{\rho}_t = \min(\rho_t, c_\rho)\) with modest \(c_\rho\) (2--5).
  \item Norm clipping of traces (or per-step projection) to keep \(\|e_t\|\) bounded.
  \item Hidden-state recomputation per-batch (re-running RNN prefix) to reduce staleness.
  \item Masking out transitions that cross episode boundaries or are padded.
\end{itemize}

\subsection{Algorithm (pseudocode)}
\begin{verbatim}
Algorithm GET-ROPR (critic update, forward-view lambda-returns):
Inputs: batch (obs [B,L,O], acts [B,L,A], rews [B,L], dones [B,L], beh_logp [B,L] or None)
1. q_vals, _ = critic(obs, acts)
2. with no-grad: next_acts = roll(acts, -1); next_q, _ = target_critic(obs, next_acts)
3. values = concat(q_vals, next_q[:, -1:])  # [B, L+1]
4. compute pi_logp = policy.log_prob(obs, acts) if policy given
5. rhos = exp(pi_logp - beh_logp) (clamp to c_rho if provided)
6. returns = lambda_returns(rews, values, dones, gamma, lam, rhos=rhos, c_rho=c_rho)
7. loss = 0.5 * (returns - q_vals).pow(2).mean()
8. optimizer.zero_grad(); loss.backward(); clip_grad_norm(); optimizer.step()
\end{verbatim}

\noindent The implementation in this repo follows the forward-view loss for
simplicity and numerical stability while maintaining conceptual equivalence
with backward-view trace accumulation.

\section{Implementation details}
We provide a compact PyTorch reference (see repository) with:
\begin{itemize}
\item GRU-based recurrent critic and RNN Gaussian policy with tanh-squash and exact log-prob correction.
\item Lambda-return loss with optional IS clipping and masking for episode ends.
\item Twin critics (TD3-style) and a minimal SAC actor update with automatic temperature (\(\alpha\)) tuning.
\end{itemize}
Key engineering choices: recompute hidden states per batch, clip traces and gradients, and use soft target updates.

\section{Experiments and Evaluation}
We evaluate GET-ROPR on two families of partially observable benchmarks:
\begin{itemize}
  \item Flicker-Atari: selected Atari games (Pong, Breakout) with intermittent frame masking (flicker) to induce partial observability.
  \item DMControl occlusion: continuous control tasks with sensor occlusion or missing observations (Reacher, Walker variants).
\end{itemize}

\subsection{Baselines and ablations}
Baselines: (i) RESeL-style recurrent off-policy (n-step TBPTT), (ii) recurrent SAC without traces, (iii) forward-view lambda-returns with \(\lambda=0\) (i.e., 1-step). Ablations: vary \(\lambda\), \(c_\rho\), trace clipping policy, and hidden recompute frequency.

\subsection{Evaluation protocol and metrics}
For each configuration we run 5 seeds. We evaluate every 10k environment steps and report:
\begin{itemize}
  \item return mean ± std, aggregated across seeds;
  \item wall-clock time per update and total GPU hours;
  \item trace norm statistics (mean, p95) and fraction of clipped steps;
  \item IS ratio mean and p95, and fraction of clipped ratios;
  \item hidden-state staleness mean/p95 (using `compute_staleness`).
\end{itemize}

\subsection{Results reporting (placeholders)}
We include tables and plotting templates to be filled after running experiments. The report contains placeholder tables and figures to be completed with CSV-derived plots (see `scripts/plot_from_csv.py` suggestion in the README). Detailed ablation tables and per-task plots are recommended for the final manuscript.

\section{Ablations and Hyperparameters}
Important ablations include \(\lambda\)-sweep, IS clip \(c_\rho\), trace clip \(c_e\), sequence length and hidden refresh toggles. Default configs are in `src/config.py`.

\section{Limitations and Ethical Considerations}
While GET-ROPR shows promise, it has limitations:
\begin{itemize}
  \item Computational overhead: traces and recomputation add wall-clock and memory cost compared to vanilla n-step implementations.
  \item Tuning sensitivity: additional hyperparameters (\(\lambda\), \(c_\rho\), trace clipping) require careful tuning per domain.
  \item Applicability: methods that rely on stored behavior log-probs require behavior logging in the replay buffer which may not always be available.
\end{itemize}

\paragraph{Ethical considerations}
Long-horizon credit assignment improves agent performance in partially observable tasks, which could be applied to both benign and potentially sensitive domains (e.g., automated systems, surveillance). Researchers should be mindful of misuse and adhere to the usual guidelines for safe and responsible deployment, including transparency, auditing, and human oversight.

\section{Conclusion}
GET-ROPR operationalizes eligibility traces for modern recurrent off-policy RL. The reference implementation, unit tests, and documentation provided with this assignment make it straightforward to reproduce and extend the experiments; remaining work focuses on large-scale ablations and reporting robust statistics across seeds.

\appendix
\section{Derivations and Proofs}
This appendix gives detailed derivations used in the body, sketches of proofs, and practical assumptions required for the method to behave stably in practice.

\subsection{Full derivation: forward-backward equivalence}
We give a step-by-step derivation showing that the gradient of the forward-view
lambda-return MSE equals the backward-view accumulation using eligibility traces.
Define the forward objective
\[
\mathcal{L} = \tfrac12 \sum_{t=0}^{L-1} (G_t^{\lambda} - Q_t)^2,
\qquad G_t^{\lambda} = Q_t + \sum_{n=1}^{L-t} (\gamma\lambda)^{n-1} \Big(\prod_{i=1}^{n-1} \bar{\rho}_{t+i}\Big) \delta_{t+n-1}.
\]
Differentiate w.r.t. parameters \(\phi\):
\[
\nabla_\phi \mathcal{L} = \sum_t (G_t^{\lambda} - Q_t) (\nabla_\phi G_t^{\lambda} - \nabla_\phi Q_t).
\]
Focus on the term involving \(\nabla_\phi G_t^{\lambda}\). Using linearity:
\[
\nabla_\phi G_t^{\lambda} = \nabla_\phi Q_t + \sum_{n=1}^{L-t} (\gamma\lambda)^{n-1} \Big(\prod_{i=1}^{n-1} \bar{\rho}_{t+i}\Big) \nabla_\phi \delta_{t+n-1}.
\]
Plugging into the gradient and rearranging sums, collect coefficients multiplying a
fixed \(\nabla_\phi Q_s\). The contribution of \(\nabla_\phi Q_s\) to \(\nabla_\phi \mathcal{L}\) arises from two sources: (i) the direct term when \(t=s\) and (ii) the terms where \(Q_s\) appears inside future returns (through \(\delta\)s). After algebraic reindexing,
one obtains
\[
\nabla_\phi \mathcal{L} = -\sum_s \Big[ \sum_{t\le s} (G_t^{\lambda} - Q_t) (\gamma\lambda)^{s-t} \prod_{i=t+1}^{s} \bar{\rho}_i \Big] \nabla_\phi Q_s.
\]
Recognizing \((G_t^{\lambda} - Q_t) = -\delta_t\) reorganizes this expression into
\(\sum_s \delta_s e_s\) where
\[
e_s = \bar{\rho}_s \big(\lambda\gamma (1-d_{s-1}) e_{s-1} + \nabla_\phi Q_s\big)
\]
as desired. The rigorous interchange of sums and the handling of clipped ratios can
be justified under mild boundedness conditions on rewards and IS ratios.

\subsection{Actor update and gradient flow}
For a stochastic policy (SAC-style), the actor objective is
\[
J_\pi(\theta) = \mathbb{E}_{o,h\sim D} \big[ \alpha \, \log \pi_\theta(a|o,h) - Q_\phi(o,h,a) \big]_{a\sim\pi_\theta},
\]
where \(D\) denotes replay sampling. The actor gradient uses \(\nabla_a Q_\phi\) which depends on the critic's parameters \(\phi\) and implicitly on the effect of traces only through the critic optimization (the traces alter \(\phi\) updates but do not appear directly in the policy gradient). In practice,
using traces yields a critic with improved long-horizon credit assignment, which
improves the fidelity of \(\nabla_a Q\) used in the actor update. For deterministic policy gradients (DDPG/TD3-style), the update becomes
\[
\nabla_\theta J \approx \mathbb{E}[\nabla_a Q_\phi(o,h,a)|_{a=\pi_\theta(o,h)} \nabla_\theta \pi_\theta(o,h)].
\]
We note that actor updates remain standard; GET-ROPR's main change is in critic targets and gradients.

\subsection{Complexity analysis}
Let \(B\) be batch size (sequences), \(L\) sequence length, and \(H\) hidden dimension.
- Time per update: forward pass O(B L H^2) (RNN+heads) ++ lambda-return processing O(B L) for returns; recompute adds roughly one additional forward pass over prefix per sequence. Overall overhead vs standard n-step: ~1.5×–2× depending on recompute frequency and implementation.
- Memory: storing intermediate values for autograd dominates standard costs; using forward-view returns avoids storing per-timestep autograd graphs for multi-step backprop, reducing memory compared to naive long-BPTT; gradient checkpointing can further reduce peak memory at the cost of extra recompute.

\subsection{Implementation checklist}
Practical tips:
\begin{itemize}
\item Vectorize sequence operations and use `torch.compile`/`functorch` where possible.
\item Compute lambda-returns in reverse order to avoid storing per-step grads.
\item Apply per-step masking for episodes and truncate traces at episode boundaries.
\item Log `trace_norm_mean`, `trace_clip_frac`, `rho_clip_frac`, and `staleness_mean` during training.
\end{itemize}

\subsection{Example command lines (local reproducibility)}
\begin{verbatim}
# Small smoke run on CPU
python train_sac_full.py --device cpu --out outputs/ca7_smoke

# Experiment run on GPU
python train_sac_full.py --device cuda --out outputs/ca7_run --seq-len 64 --batch-size 64
\end{verbatim}



\subsection{IS clipping and bias bound}
Clipping importance ratios at \(c_\rho\) introduces bias. For sequences of length at most \(L\), a simple bound on the forward-view bias (difference between clipped and unclipped lambda-returns) can be written as
\[
|G_t^{\lambda,\mathrm{clipped}} - G_t^{\lambda}| \le \frac{\gamma^{L-t} R_{\max}}{1-\gamma} \big(1 - c_\rho^{L-t}\big),
\]
which follows from bounding the effect of replacing a product of ratios by clipped versions and using \(|\delta|\le R_{\max}/(1-\gamma)\). In practice moderate values of \(c_\rho\) (2--5) reduce variance substantially with small empirical bias.

\subsection{Trace-norm bound}
Assume \(\|\nabla_\phi Q_t\|\le G_{\max}\) and \(\bar{\rho}_t\le c_\rho\). Then
\[
\|e_t\| \le c_\rho \sum_{k=0}^t (\lambda\gamma c_\rho)^k G_{\max} = \frac{c_\rho G_{\max}}{1-\lambda\gamma c_\rho},
\]
which requires \(\lambda\gamma c_\rho < 1\) for finite bound. This motivates conservative choices of \(\lambda\) and \(c_\rho\) if \(\gamma\) is near 1.

\section{Algorithmic Details and Pseudocode}
Below is a concise pseudocode reference for the critic update using lambda-returns (forward-view) and hidden-state recomputation. This is intentionally simple and mirrors the implementation in the provided reference code.
\begin{verbatim}
# Inputs: batch of sequences (obs [B,L,O], acts [B,L,A], rews [B,L], dones [B,L])
# critic, target_critic, policy, optimizer, gamma, lam, c_rho
q_vals, _ = critic(obs, acts)  # [B,L]
with torch.no_grad():
    next_acts = roll(acts, -1)
    next_q, _ = target_critic(obs, next_acts)
    values = concat(q_vals, next_q[:, -1:])  # [B, L+1]
# compute pi logp for IS ratios if behavior logps stored
rhos = compute_rhos(policy, obs, acts, beh_logp, c_rho)
returns = lambda_returns(rews, values, dones, gamma, lam, rhos=rhos, c_rho=c_rho)
loss = 0.5 * (returns - q_vals).pow(2).mean()
optimizer.zero_grad(); loss.backward(); clip_grad_norm(critic, max_norm); optimizer.step()
\end{verbatim}

\section{Experimental Protocol}
We recommend the following reproducible protocol for running the reported experiments:

\begin{enumerate}
  \item Use fixed seeds (PyTorch, NumPy, envs) and record the git commit hash.
  \item Train on at least 3--5 random seeds per condition (5 recommended for paper-quality results).
  \item Evaluate policy every N steps (e.g., 10k steps) by running 10 deterministic evaluation episodes and report mean ± std across seeds.
  \item Log TD-errors, trace norms, IS ratios, hidden-state staleness and standard learning metrics (returns, losses, grad norms).
  \item Save raw CSV logs and a config snapshot with each checkpoint.
\end{enumerate}

\section{Hyperparameters (Suggested)}
\begin{tabular}{l l}
Parameter & Recommended value / range \\
\hline
Sequence length $L$ & 32--64 \\
$\lambda$ & 0.8--0.95 \\
$\gamma$ & 0.99 (Atari), 0.995 (DMControl) \\
c_{\rho} (IS clip) & 2--5 \\
$c_e$ (trace clip) & 1--5 (norm clip) \\
Batch size (sequences) & 32--128 \\
Learning rate (actor/critic) & 3e-4 (tune) \\
Averaging seeds & 3--10 (paper: 5) \\
\end{tabular}

\section{Reproducibility and Compute}
All code, demo scripts and configs are included in this repository. To reproduce experiments:
\begin{enumerate}
  \item Create environment and install requirements (see `requirements.txt`).
  \item Run training script `python train_sac_full.py --device cuda --out outputs/ca7` with a desired config.
  \item Use `scripts/plot_from_csv.py` (suggested) to produce figures from `outputs/ca7/train_log.csv`.
\end{enumerate}

Hardware notes: experiments were designed to run on a single GPU (A100 or similar); typical small runs (seq_len=32, batch=64, hidden=256) complete 100k updates in several hours on A100; scale accordingly.

\section{Additional Figures and Templates}
Placeholders for figures and result tables are included in the main manuscript; after running experiments, populate `pictures/` with `fig_returns.png`, `fig_metrics.png`, and `fig_trace_heatmap.png`.

\section{Full Checklist}
- [ ] Run full ablation suite (lambda, c_rho, seq_len, refresh) across 3--5 seeds each.
- [ ] Generate plots and fill result tables in `report.tex`.
- [ ] Archive logs, checkpoints, and config snapshots for reproducibility.

\section{Code provenance}
Implementation files of note:
\begin{itemize}
  \item `src/model.py` — recurrent critic and actor models
  \item `src/losses.py` — lambda-return and critic loss
  \item `src/data.py` — sequence buffer utilities
  \item `train_sac_full.py` — demo runner and logging
  \item `tests/` — unit and smoke tests
\end{itemize}

\section{Acknowledgements}
We thank the class instructors and peers for feedback and the authors whose work inspired this project.



\paragraph{Code and tests} The implementation includes unit smoke tests in `tests/` (run via `pytest -q`). The demo training script `train_sac_full.py` logs `returns_mean` to `outputs/ca7/train_log.csv` and checkpoints models under `outputs/ca7/`. Figures and tables in this manuscript are placeholders until experiment runs are executed and saved to `pictures/`.

\bibliographystyle{plain}
\bibliography{refs}
\end{document}

















